\subsection{半单模中的重数与长度}

命题 10. —— 设 M 是一个半单 A-模。设 \((\mathrm{M}_{i})_{i\in\mathrm{I}}\) 是以 M 为直和的单子模族。下列性质等价：

\begin{enumerate}
\def\labelenumi{(\roman{enumi})}
\item
  M 的长度有限。
\item
  M 是阿廷的。
\item
  M 是诺特的。
\item
  M 是有限生成的。
\item
  I 是有限的。
\end{enumerate}

若 M 具有这些性质，则 M 的长度等于 I 的基数。

若集合 I 有限，则 M 具有性质 (i)、(ii)、(iii)、(iv)。设集合 I 无限。由 VIII, 第 2 页的例 2，模 M 既非阿廷的亦非诺特的；由于每个有限长度的模都是阿廷的且诺特的（VIII, 第 2 页，命题 1），M 也不是有限长度的。最后，M 的每个元素都属于有限个子模 \(M_{i}\) 之和，故 M 不是有限生成的。这就证明了性质 (i) 至 (v) 的等价性。若这些性质成立，则我们有 \(\text{long}(M) = \sum_{i \in I} \text{long}(M_{i}) = \text{Card}(I)\)（II, §1, No.~10, 第 213 页，推论 5）。

命题 11. —— 设 M 是一个半单 A-模，它是单子模族 \((\mathrm{M}_{i})_{i\in\mathrm{I}}\) 的直和。对每个 \(\lambda\in S\)，用 \(\mathrm{I}(\lambda)\) 表示这样的指标 \(i\in I\) 组成的集合：\(M_{i}\) 属于类 \(\lambda\)。\(\mathrm{I}(\lambda)\) 的基数等于左 \(D_{\lambda}\)-向量空间 \(\mathrm{Hom}_{\mathrm{A}}(\mathrm{S}_{\lambda},\mathrm{M})\) 的维数。

A 的 \(\lambda\) 型同型分量同构于 \(\mathrm{S}_{\lambda}^{(\mathrm{I}(\lambda))}\)（VIII, 第 65 页，命题 4, b)）。\(D_{\lambda}\)-向量空间 \(\mathrm{Hom}_{\mathrm{A}}(\mathrm{S}_{\lambda},\mathrm{M})\) 可与 \(\mathrm{Hom}_{\mathrm{A}}(\mathrm{S}_{\lambda},\mathrm{M}_{\lambda})\) 等同，故同构于 \(\mathrm{D}_{\lambda}^{(\mathrm{I}(\lambda))}\)（No.~2）。这就证明了本命题。

每个单模都是原始模（VIII, 第 45 页），故每个半单模都是半原始模。设 M 是一个半单 A-模，\(\lambda \in \mathcal{S}\)。\(\lambda\) 在 M 中的重数是 VIII, 第 34 页定义的 \([\mathrm{M}:\lambda]\)，称为 \(\lambda\) 在 M 中的原始重数。命题 11 可以转述为等式

\[
[ \mathrm{M}: \lambda ] = \dim_ {\mathrm{D} _ {\lambda}} (\operatorname{Hom} _ {\mathrm{A}} (\mathrm{S} _ {\lambda}, \mathrm{M})).\tag{11}
\]

更一般地，若 \(((\mathrm{V}_{\lambda})_{\lambda \in \mathcal{S}},\alpha)\) 是 M 的一个描述，则 \([\mathrm{M}:\lambda ]\) 等于 \(\dim_{\mathrm{D}_{\lambda}}(\mathrm{V}_{\lambda})\)。由 VIII, 第 68 页的命题 6，我们还有

\[
[ \mathrm{M}: \lambda ] = \dim_ {\mathrm{D} _ {\lambda}} (\operatorname{Hom} _ {\mathrm{A}} (\mathrm{M}, \mathrm{S} _ {\lambda}))\tag{12}
\]

其中重数 \([M : \lambda]\) 有限。半单 A-模 M 与 \(M'\) 同构当且仅当我们有 \([M : \lambda] = [M' : \lambda]\)（对每个 \(\lambda \in S\)）。

设 M 是一个半单 A-模。存在一个具有下列性质的基数 I：对 M 分解为单模直和的每个分解 \(M = \bigoplus_{i \in I} M_i\)，I 的基数都等于 I（VIII, 第 34 页，推论 2）。这个基数称为半单 A-模 M 的长度，记作 \(\operatorname{long}_{\mathrm{A}}(\mathrm{M})\) 或 \(\operatorname{long}(\mathrm{M})\)。当 M 的长度有限时，由命题 10，这个定义与 II, §1, No.~10, 第 212 页中的定义相容。

单 A-模是长度为 1 的半单 A-模，并且我们有公式

\[
\operatorname{long} _ {\mathrm{A}} \left(\bigoplus_ {j \in \mathrm{J}} \mathrm{M} _ {j}\right) = \sum_ {j \in \mathrm{J}} \operatorname{long} _ {\mathrm{A}} (\mathrm{M} _ {j})\tag{13}
\]

其中 \((\mathrm{M}_j)_{j\in \mathrm{J}}\) 是任意一族半单 A-模。由命题 11，我们有

\[
\mathrm{long} _ {\mathrm{A}} (\mathrm{M}) = \sum_ {\lambda \in \mathcal {S}} \dim_ {\mathrm{D} _ {\lambda}} \mathrm{Hom} _ {\mathrm{A}} (\mathrm{S} _ {\lambda}, \mathrm{M}).\tag{14}
\]

把这个公式应用于 \(M_{\lambda}\)，我们得到 \([M : \lambda] = \text{long}_{A}(M_{\lambda})\)（对每个 \(\lambda \in S\)）。

当 A 是一个体时，单 A-模是维数为 1 的向量空间；这时每个 A-模都是半单的（II, §7, No.~1, 第 292 页，定理 1），而它的长度就是它作为 A 上向量空间的维数（II, §7, No.~2, 第 293 页）。

